%2multibyte Version: 5.50.0.2953 CodePage: 1252

\documentclass[notes=show]{beamer}
%%%%%%%%%%%%%%%%%%%%%%%%%%%%%%%%%%%%%%%%%%%%%%%%%%%%%%%%%%%%%%%%%%%%%%%%%%%%%%%%%%%%%%%%%%%%%%%%%%%%%%%%%%%%%%%%%%%%%%%%%%%%%%%%%%%%%%%%%%%%%%%%%%%%%%%%%%%%%%%%%%%%%%%%%%%%%%%%%%%%%%%%%%%%%%%%%%%%%%%%%%%%%%%%%%%%%%%%%%%%%%%%%%%%%%%%%%%%%%%%%%%%%%%%%%%%
\usepackage{amsfonts}
\usepackage{graphicx}
\usepackage{amssymb}
\usepackage{amsmath}
\usepackage{mathpazo}
\usepackage{hyperref}
\usepackage{multimedia}

\setcounter{MaxMatrixCols}{10}
%TCIDATA{OutputFilter=LATEX.DLL}
%TCIDATA{Version=5.50.0.2953}
%TCIDATA{Codepage=1252}
%TCIDATA{<META NAME="SaveForMode" CONTENT="2">}
%TCIDATA{BibliographyScheme=Manual}
%TCIDATA{Created=Wednesday, May 01, 2013 14:54:39}
%TCIDATA{LastRevised=Wednesday, November 17, 2021 14:23:58}
%TCIDATA{<META NAME="GraphicsSave" CONTENT="32">}
%TCIDATA{<META NAME="DocumentShell" CONTENT="Other Documents\SW\Slides - Beamer">}
%TCIDATA{Language=American English}
%TCIDATA{CSTFile=beamer.cst}
%TCIDATA{PageSetup=72,72,72,83,0}
%TCIDATA{Counters=arabic,1}
%TCIDATA{AllPages=
%H=36
%F=36
%}


\newenvironment{stepenumerate}{\begin{enumerate}[<+->]}{\end{enumerate}}
\newenvironment{stepitemize}{\begin{itemize}[<+->]}{\end{itemize} }
\newenvironment{stepenumeratewithalert}{\begin{enumerate}[<+-| alert@+>]}{\end{enumerate}}
\newenvironment{stepitemizewithalert}{\begin{itemize}[<+-| alert@+>]}{\end{itemize} }
\input{tcilatex}
\usetheme{JuanLesPins}
\usecolortheme{whale}

\begin{document}

\title{Econometrics Lecture 6: \\
Endogeneity, Instrumental Variables and Two Stage Least Squares}
\author{Katerina Petrova \ \ \ \ \ \ \ \ \ \ \ \ \ \ \ \ \ \ \ \ \ \ \ \ \ \
\ \ \ \ \ \ \ \ \ \ \ \ \ \ \ \ \ \ \ \ \ \ \ \ \ \ \ \ \ \ \ \ \ \ \ \ \ \
\ \ \ \ \ \ \ \ \ \ \ \ \ \ \ \ \ \ \ \ \ \ \ }
\date{}
\maketitle

\section{The Model}

%TCIMACRO{\TeXButton{BeginFrame}{\begin{frame}}}%
%BeginExpansion
\begin{frame}%
%EndExpansion

\QTR{frametitle}{Classical Assumptions for OLS}

\begin{itemize}
\item We have been working with the following model 
\begin{equation}
y=X\beta +\varepsilon  \label{model}
\end{equation}

\item Let's recall the assumptions we made
\end{itemize}

\begin{enumerate}
\item The model in (\ref{model}) is correctly specified, and crucially,
linear in $\beta $ and $\varepsilon .$

\item The columns of $X$ are linearly independent, so that $rank(X^{\prime
}X)=rank(X)=k<n.$

\item $\varepsilon _{i}$ are i.i.d. variables and $X$ are exogenous
identically distributed regressors, such that $\mathbb{E}[\varepsilon
_{i}|X]=0$ and $\mathbb{E}[\varepsilon _{i}^{2}|X]=\sigma ^{2}.$

\item $p\lim_{n\rightarrow \infty }(\frac{1}{n}X^{\prime
}X)=p\lim_{n\rightarrow \infty }\frac{1}{n}\sum_{i=1}^{n}\mathbf{x}_{i}%
\mathbf{x}_{i}^{\prime }=\mathbb{E[}\mathbf{x}_{1}\mathbf{x}_{1}^{\prime
}]=Q,$ where $\left\Vert Q\right\Vert <\infty $ and $Q$ is positive definite
matrix.
\end{enumerate}

%TCIMACRO{\TeXButton{EndFrame}{\end{frame}}}%
%BeginExpansion
\end{frame}%
%EndExpansion

\section{Endogeneity}

%TCIMACRO{\TeXButton{BeginFrame}{\begin{frame}}}%
%BeginExpansion
\begin{frame}%
%EndExpansion

\QTR{frametitle}{Relaxing Assumption 3}

\begin{itemize}
\item In this lecture, we will maintain all the classical assumptions above
but we will relax the assumption of exogenous regressors

\item In particular, we relax the assumption that $X$ are exogenous
regressors, that is $\mathbb{E}[\varepsilon |X]\neq 0.$
\end{itemize}

%TCIMACRO{\TeXButton{EndFrame}{\end{frame}}}%
%BeginExpansion
\end{frame}%
%EndExpansion

%TCIMACRO{\TeXButton{BeginFrame}{\begin{frame}}}%
%BeginExpansion
\begin{frame}%
%EndExpansion

\QTR{frametitle}{What has changed?}

$\widehat{\beta }_{n}^{OLS}$ is now biased.

\begin{eqnarray*}
\mathbb{E}[\widehat{\beta }_{n}^{OLS}] &=&\beta +\mathbb{E}[(X^{\prime
}X)^{-1}X^{\prime }\varepsilon ] \\
&=&\beta +\mathbb{E[E[}(X^{\prime }X)^{-1}X^{\prime }\varepsilon |X]] \\
&=&\beta +\mathbb{E[}(X^{\prime }X)^{-1}X^{\prime }\underset{\neq 0}{%
\underbrace{\mathbb{E[}\varepsilon |X]}]}\neq \beta
\end{eqnarray*}

%TCIMACRO{\TeXButton{EndFrame}{\end{frame}}}%
%BeginExpansion
\end{frame}%
%EndExpansion

%TCIMACRO{\TeXButton{BeginFrame}{\begin{frame}}}%
%BeginExpansion
\begin{frame}%
%EndExpansion

\QTR{frametitle}{What has changed?}

\begin{itemize}
\item Crucially, it is no longer consistent!
\end{itemize}

$p\lim_{n\rightarrow \infty }\widehat{\beta }_{n}^{OLS}=\beta
+p\lim_{n\rightarrow \infty }\left[ \left( \frac{1}{n}\tsum%
\nolimits_{i=1}^{n}x_{i}x_{i}^{\prime }\right) ^{-1}\frac{1}{n}%
\tsum\nolimits_{i=1}^{n}x_{i}\varepsilon _{i}\right] $%
\begin{eqnarray*}
&=&\beta +p\lim_{n\rightarrow \infty }\left( \frac{1}{n}\tsum%
\nolimits_{i=1}^{n}x_{i}x_{i}^{\prime }\right) ^{-1}p\lim_{n\rightarrow
\infty }\left( \frac{1}{n}\tsum\nolimits_{i=1}^{n}x_{i}\varepsilon
_{i}\right) \\
&&\overset{\text{WLLN}}{=}\beta +Q^{-1}\underset{\neq 0}{\underbrace{\mathbb{%
E[}x_{i}\varepsilon _{i}]}} \\
&\neq &\beta
\end{eqnarray*}

%TCIMACRO{\TeXButton{EndFrame}{\end{frame}}}%
%BeginExpansion
\end{frame}%
%EndExpansion

%TCIMACRO{\TeXButton{BeginFrame}{\begin{frame}}}%
%BeginExpansion
\begin{frame}%
%EndExpansion

\QTR{frametitle}{Consequences}

\begin{itemize}
\item Endogeneity of the regressors $X$ invalidates the OLS procedure even
asymptotically

\item Sources: measurement error in $X,$ simultaneous equation models,
unobserved causal variable correlated with $y$ and $X$

\item Solution: turns out that if we can find a variable $Z$ (of the same
dimension as $X$), such that $Z$ satisfies the original exogeneity
assumption $\mathbb{E}[\varepsilon _{i}|z_{i}]=0,$ and in the same time $%
\mathbb{E[}z_{i}x_{i}^{\prime }]\neq 0$ (i.e. $Z$ is correlated with $X$),
then we can construct a new estimator, which will be unbiased and consistent

\item In this case, we will refer to $Z$ as instrumental variable, we will
refer to condition $\mathbb{E}[\varepsilon _{i}|z_{i}]=0$ as a validity of
the instrument and to condition $\mathbb{E[}z_{i}x_{i}^{\prime }]\neq 0$ as
a relevance of the instrument

\item The resulting estimator is called Instrumental variables (or IV)
estimator
\end{itemize}

%TCIMACRO{\TeXButton{EndFrame}{\end{frame}}}%
%BeginExpansion
\end{frame}%
%EndExpansion

\section{The IV\ estimator}

%TCIMACRO{\TeXButton{BeginFrame}{\begin{frame}}}%
%BeginExpansion
\begin{frame}%
%EndExpansion

\QTR{frametitle}{New assumptions for IV}

\begin{itemize}
\item So let's modify the assumptions

\item Assumption 3$^{\text{*}}$: $\varepsilon _{i}$ are random variables,
such that $\mathbb{E}[\varepsilon _{i}|x_{i}]\neq 0$, but we can find a full
rank matrix of identically distributed instruments $Z=\left[ 
\begin{array}{c}
z_{1}^{\prime } \\ 
\vdots \\ 
z_{n}^{\prime }%
\end{array}%
\right] ,$ where $z_{i}$ is of the same size ($k\times 1$) as $x_{i},$ are
exogenous regressors$,$ such that $\mathbb{E}[\varepsilon _{i}|z_{i}]=0$ and 
$\mathbb{E}[\varepsilon _{i}^{2}|z_{i}]=\sigma ^{2}.$

\item Assumption 4$^{\text{*}}$: $p\lim_{n\rightarrow \infty }(\frac{1}{n}%
Z^{\prime }X)=p\lim_{n\rightarrow \infty }\frac{1}{n}%
\sum_{i=1}^{n}z_{i}x_{i}^{\prime }=S,$ where $\left\Vert S\right\Vert
<\infty $ and $rank(S)=k.$
\end{itemize}

%TCIMACRO{\TeXButton{EndFrame}{\end{frame}}}%
%BeginExpansion
\end{frame}%
%EndExpansion

%TCIMACRO{\TeXButton{BeginFrame}{\begin{frame}}}%
%BeginExpansion
\begin{frame}%
%EndExpansion

\QTR{frametitle}{Remarks}

\begin{itemize}
\item If $z_{i}$ is of dimension smaller than $x_{i},$ the IV estimator will
not be defined and we will say that the model is under-identified

\item For the IV to work, we need to assume that $z_{i}$ is of the same
dimension of $x_{i}$; in this case, we say that the model is just-identified

\item On the other hand, if we had more valid and relevant instruments than
endogenous variables, then, we will arrive at the Two Stage Least Squares
(2SLS) estimator, which is an over-identified model and will turn out to be
efficient

\item If only one of the regressors in $X$ was endogenous, we could
partition $X$ and instrument only the endogenous part
\end{itemize}

%TCIMACRO{\TeXButton{EndFrame}{\end{frame}}}%
%BeginExpansion
\end{frame}%
%EndExpansion

%TCIMACRO{\TeXButton{BeginFrame}{\begin{frame}}}%
%BeginExpansion
\begin{frame}%
%EndExpansion

\QTR{frametitle}{The IV estimator}

\begin{itemize}
\item The IV estimator solves the following moment condition%
\begin{gather*}
\mathbb{E}[z_{i}\varepsilon _{i}]=0 \\
\mathbb{E}[z_{i}(y_{i}-x_{i}^{\prime }\beta )]=0 \\
\mathbb{E}[z_{i}x_{i}^{\prime }]\beta =\mathbb{E}[z_{i}y_{i}] \\
\widehat{\beta }^{IV}=\left( \frac{1}{n}\sum_{i=1}^{n}z_{i}x_{i}^{\prime
}\right) ^{-1}\frac{1}{n}\sum_{i=1}^{n}z_{i}y_{i}
\end{gather*}

\item We can invert the $\mathbb{E}[z_{i}x_{i}^{\prime }]$ matrix, because
of the additional assumption we made in 4$^{\text{*}}$.

\item Obviously, we don't know the population quantities $\mathbb{E}%
[z_{i}x_{i}^{\prime }]$ and $\mathbb{E}[z_{i}y_{i}],$ but we can replace
them with their sample counterparts: $\frac{1}{n}\sum_{i=1}^{n}z_{i}x_{i}^{%
\prime }$ and $\frac{1}{n}\sum_{i=1}^{n}z_{i}y_{i}.$

$\widehat{\beta }^{IV}=\left( \frac{1}{n}\sum_{i=1}^{n}z_{i}x_{i}^{\prime
}\right) ^{-1}\frac{1}{n}\sum_{i=1}^{n}z_{i}y_{i}=\left( Z^{\prime }X\right)
^{-1}Z^{\prime }Y.$
\end{itemize}

%TCIMACRO{\TeXButton{EndFrame}{\end{frame}}}%
%BeginExpansion
\end{frame}%
%EndExpansion

\section{Large Sample Properties}

%TCIMACRO{\TeXButton{BeginFrame}{\begin{frame}}}%
%BeginExpansion
\begin{frame}%
%EndExpansion

\QTR{frametitle}{IV: Asymptotic Properties}

\begin{gather*}
p\lim_{n\rightarrow \infty }\widehat{\beta }^{IV}=\beta +p\lim_{n\rightarrow
\infty }\left[ \left( \frac{1}{n}\tsum\nolimits_{i=1}^{n}z_{i}x_{i}^{\prime
}\right) ^{-1}\frac{1}{n}\tsum\nolimits_{i=1}^{n}z_{i}\varepsilon _{i}\right]
\\
=\beta +\left( p\lim_{n\rightarrow \infty }\left( \frac{1}{n}%
\tsum\nolimits_{i=1}^{n}z_{i}x_{i}^{\prime }\right) \right)
^{-1}p\lim_{n\rightarrow \infty }\left( \frac{1}{n}\tsum%
\nolimits_{i=1}^{n}z_{i}\varepsilon _{i}\right) \text{ by CMT} \\
=\beta +\underset{=S^{-1}}{\underbrace{\mathbb{E}[z_{1}x_{1}^{\prime }]^{-1}}%
}\underset{\mathbb{E[}z_{i}\mathbb{E}[\varepsilon _{i}|z_{i}]]=0\text{ by 3}%
^{\text{*}}}{\underbrace{\mathbb{E[}z_{i}\varepsilon _{i}]}} \\
=\beta
\end{gather*}

%TCIMACRO{\TeXButton{EndFrame}{\end{frame}}}%
%BeginExpansion
\end{frame}%
%EndExpansion

%TCIMACRO{\TeXButton{BeginFrame}{\begin{frame}}}%
%BeginExpansion
\begin{frame}%
%EndExpansion

\QTR{frametitle}{IV: Asymptotic Properties}

\begin{eqnarray*}
\sqrt{n}(\widehat{\beta }^{IV}-\beta ) &=&\left( \frac{1}{n}%
\tsum\nolimits_{i=1}^{n}z_{i}x_{i}^{\prime }\right) ^{-1}\frac{1}{\sqrt{n}}%
\tsum\nolimits_{i=1}^{n}z_{i}\varepsilon _{i} \\
\left( \frac{1}{n}\tsum\nolimits_{i=1}^{n}z_{i}x_{i}^{\prime }\right) ^{-1}
&\rightarrow &_{p}\text{ }S^{-1}\text{ by 4}^{\text{*}}\text{ and CMT} \\
\frac{1}{\sqrt{n}}\tsum\nolimits_{i=1}^{n}z_{i}\varepsilon _{i} &\rightarrow
&_{d}\text{ }\mathcal{N}(0,var(z_{1}\varepsilon _{1})) \\
&\rightarrow &_{d}\text{ }\mathcal{N}(0,\sigma ^{2}Q_{z}) \\
\sqrt{n}(\widehat{\beta }^{IV}-\beta ) &\rightarrow &_{d}\mathcal{N}%
(0,\sigma ^{2}S^{-1}Q_{z}\left( S^{\prime }\right) ^{-1})\text{ by Slutsky }
\end{eqnarray*}%
with $Q_{z}=\mathbb{E[}z_{i}z_{i}^{\prime }].$

%TCIMACRO{\TeXButton{EndFrame}{\end{frame}}}%
%BeginExpansion
\end{frame}%
%EndExpansion

%TCIMACRO{\TeXButton{BeginFrame}{\begin{frame}}}%
%BeginExpansion
\begin{frame}%
%EndExpansion

\QTR{frametitle}{IV: Discussion}

\begin{itemize}
\item We saw that under endogeneity, OLS is invalid even asymptotically

\item However, given the right instruments (valid and relevant), we are back
in business, and $\widehat{\beta }^{IV}$ is consistent, and we can use the
asymptotic distribution to do testing or construct confidence intervals

\item Everything we have done on testing with OLS applies immediately, by
replacing $\widehat{\beta }^{OLS}\ $with $\widehat{\beta }^{IV}$ and the
standard error of $\widehat{\beta }^{OLS}$ with the standard error of $%
\widehat{\beta }^{IV}.$

\item What about when we have more instruments that regressors in $X,$ i.e. $%
Z$ is a $n\times l$ matrix, where $l>k?$ Can we gain efficiency by using all
of them?

\item The IV procedure is not directly applicable, since $(Z^{\prime }X)$ is
a $l\times k$ matrix, so we cannot invert it

\item We could, however, do two stage least squares (2SLS)
\end{itemize}

%TCIMACRO{\TeXButton{EndFrame}{\end{frame}}}%
%BeginExpansion
\end{frame}%
%EndExpansion

\section{2SLS}

%TCIMACRO{\TeXButton{BeginFrame}{\begin{frame}}}%
%BeginExpansion
\begin{frame}%
%EndExpansion

\QTR{frametitle}{2SLS}

\begin{itemize}
\item Stage one: regress $Z$ on $X$ using OLS%
\begin{eqnarray*}
\underset{n\times k}{X} &=&\underset{n\times l}{Z}\text{ }\underset{l\times k%
}{\gamma }+\underset{n\times k}{v} \\
\widehat{\gamma } &=&(Z^{\prime }Z)^{-1}Z^{\prime }X\text{ \ }\Rightarrow 
\widehat{X}=Z\widehat{\gamma }
\end{eqnarray*}

\item Stage 2: use the prediction $\widehat{X}$ instead of the endogenous $X$
in the model and do OLS%
\begin{eqnarray*}
Y &=&\widehat{X}\beta +\varepsilon \text{ \ \ }\Rightarrow \widehat{\beta }%
^{2SLS}=(\widehat{X}^{\prime }\widehat{X})^{-1}\widehat{X}^{\prime }y \\
&=&(\widehat{\gamma }^{\prime }Z^{\prime }Z\widehat{\gamma })^{-1}\widehat{%
\gamma }^{\prime }Z^{\prime }y \\
&=&(X^{\prime }Z(Z^{\prime }Z)^{-1}Z^{\prime }Z(Z^{\prime }Z)^{-1}Z^{\prime
}X)^{-1}X^{\prime }Z(Z^{\prime }Z)^{-1}Z^{\prime }y \\
&=&(X^{\prime }\underset{P_{Z}}{\underbrace{Z(Z^{\prime }Z)^{-1}Z^{\prime }}}%
X)^{-1}X^{\prime }\underset{P_{Z}}{\underbrace{Z(Z^{\prime }Z)^{-1}Z^{\prime
}}}y \\
&\Rightarrow &\widehat{\beta }^{2SLS}=(X^{\prime }P_{Z}X)^{-1}X^{\prime
}P_{Z}y
\end{eqnarray*}
\end{itemize}

%TCIMACRO{\TeXButton{EndFrame}{\end{frame}}}%
%BeginExpansion
\end{frame}%
%EndExpansion

%TCIMACRO{\TeXButton{BeginFrame}{\begin{frame}}}%
%BeginExpansion
\begin{frame}%
%EndExpansion

\QTR{frametitle}{2SLS}

\begin{itemize}
\item What happens with $\widehat{\beta }^{2SLS}$ when $l=k?$

\item In that case $X^{\prime }Z$ is a square full rank matrix, so that we
can open the bracket%
\begin{eqnarray*}
\widehat{\beta }^{2SLS} &=&(X^{\prime }Z(Z^{\prime }Z)^{-1}Z^{\prime
}X)^{-1}X^{\prime }Z(Z^{\prime }Z)^{-1}Z^{\prime }y \\
&=&(Z^{\prime }X)^{-1}(Z^{\prime }Z)(X^{\prime }Z)^{-1}X^{\prime
}Z(Z^{\prime }Z)^{-1}Z^{\prime }y \\
&=&(Z^{\prime }X)^{-1}Z^{\prime }y=\widehat{\beta }^{IV}
\end{eqnarray*}

\item 2SLS is a special case of optimal GMM.
\end{itemize}

%TCIMACRO{\TeXButton{EndFrame}{\end{frame}}}%
%BeginExpansion
\end{frame}%
%EndExpansion

%TCIMACRO{\TeXButton{BeginFrame}{\begin{frame}}}%
%BeginExpansion
\begin{frame}%
%EndExpansion

\QTR{frametitle}{2SLS Consistency}

\begin{itemize}
\item Consistency of 2SLS follows immediately since $\widehat{\beta }%
^{2SLS}=\beta +(X^{\prime }Z(Z^{\prime }Z)^{-1}Z^{\prime }X)^{-1}X^{\prime
}Z(Z^{\prime }Z)^{-1}Z^{\prime }\varepsilon $

\item $p\lim_{n\rightarrow \infty }\widehat{\beta }^{2SLS}$%
\begin{gather*}
=\beta +p\lim_{n\rightarrow \infty }\left[ \left( \frac{\tsum%
\nolimits_{i=1}^{n}x_{i}z_{i}^{\prime }}{n}\left( \frac{\tsum%
\nolimits_{i=1}^{n}z_{i}z_{i}^{\prime }}{n}\right) ^{-1}\frac{%
\tsum\nolimits_{i=1}^{n}z_{i}x_{i}^{\prime }}{n}\right) ^{-1}\right] \\
\times p\lim_{n\rightarrow \infty }\left[ \frac{\tsum%
\nolimits_{i=1}^{n}x_{i}z_{i}^{\prime }}{n}\left( \frac{\tsum%
\nolimits_{i=1}^{n}z_{i}z_{i}^{\prime }}{n}\right) ^{-1}\frac{%
\tsum\nolimits_{i=1}^{n}z_{i}\varepsilon _{i}}{n}\right] \\
=\beta +\left( S^{\prime }Q_{z}^{-1}S\right) ^{-1}S^{\prime }Q_{z}^{-1}%
\mathbb{E[}z_{i}\varepsilon _{i}]\text{ by CMT, WLLN} \\
=\beta \text{ by Assumption 3}^{\text{*}}.
\end{gather*}
\end{itemize}

%TCIMACRO{\TeXButton{EndFrame}{\end{frame}}}%
%BeginExpansion
\end{frame}%
%EndExpansion

%TCIMACRO{\TeXButton{BeginFrame}{\begin{frame}}}%
%BeginExpansion
\begin{frame}%
%EndExpansion

\QTR{frametitle}{2SLS Asymptotic Distribution}%
\begin{gather*}
\sqrt{n}\left( \widehat{\beta }^{2SLS}-\beta \right) =\left( \frac{%
\tsum\nolimits_{i=1}^{n}x_{i}z_{i}^{\prime }}{n}\left( \frac{%
\tsum\nolimits_{i=1}^{n}z_{i}z_{i}^{\prime }}{n}\right) ^{-1}\frac{%
\tsum\nolimits_{i=1}^{n}z_{i}x_{i}^{\prime }}{n}\right) ^{-1} \\
\times \frac{\tsum\nolimits_{i=1}^{n}x_{i}z_{i}^{\prime }}{n}\left( \frac{%
\tsum\nolimits_{i=1}^{n}z_{i}z_{i}^{\prime }}{n}\right) ^{-1}\frac{%
\tsum\nolimits_{i=1}^{n}z_{i}\varepsilon _{i}}{\sqrt{n}} \\
\rightarrow _{d}\left( S^{\prime }Q_{z}^{-1}S\right) ^{-1}S^{\prime
}Q_{z}^{-1}\mathcal{N}\left( 0,\sigma ^{2}Q_{z}\right) \text{ } \\
\text{{\small by CMT, Slutsky, WLLN, CLT, Ass. 3 and 4}} \\
\rightarrow _{d}\mathcal{N}\left( 0,\sigma ^{2}\left( S^{\prime
}Q_{z}^{-1}S\right) ^{-1}\right) .
\end{gather*}

%TCIMACRO{\TeXButton{EndFrame}{\end{frame}}}%
%BeginExpansion
\end{frame}%
%EndExpansion

\section{Reading}

%TCIMACRO{\TeXButton{BeginFrame}{\begin{frame}}}%
%BeginExpansion
\begin{frame}%
%EndExpansion

\QTR{frametitle}{Reading}

\begin{itemize}
\item Relevant chapters from the textbook: 12
\end{itemize}

%TCIMACRO{\TeXButton{EndFrame}{\end{frame}}}%
%BeginExpansion
\end{frame}%
%EndExpansion

\end{document}
